% Subsection 3.3.2: 标签体系

少样本检索（\texttt{CoT-FS} 推理与 \texttt{LoRA\_cot\_fs} 训练）需在题库中按算法范式找到「真正同类」的示例。\texttt{Array}、\texttt{String}、\texttt{Math} 等通用标签在题库中近似全命中，会稀释相似度信号。本文显式将标签词表切分为强标签与弱标签（见表~\ref{tab:ch3-tags}）：强标签刻画高判别力的算法骨架，弱标签保留通用类型与宽泛范式。教师模型 \texttt{deepseek-chat} 按此协议产出每条样本的 \texttt{tags}，并施加硬约束「至少含一个强标签」，否则重采样直至满足。

最终落盘 \texttt{KodCode\_annotate.jsonl}，仅在原字段基础上附加 \texttt{tags}。检索时与 \texttt{STRONG\_TAGS} 常量取交集即时派生 \texttt{strong\_tags}：训练侧 \texttt{LoRA\_cot\_fs} 与评测侧 \texttt{CoT-FS} 共享同一索引，避免示例挑选口径漂移。

\begin{table}[!htb]
  \centering
  \caption{强/弱标签词表（节选自 \texttt{annotate\_algorithm\_tags.py}）}
  \label{tab:ch3-tags}
  \footnotesize
  \hspace{-2cm}
  \begin{tabular}{p{3cm} p{13cm}}
    \toprule
    \textbf{类别} & \textbf{标签}\\
    \midrule
    核心范式（强）    & Dynamic Programming, Greedy, Backtracking, Divide and Conquer, Memoization\\
    搜索/图算法（强） & BFS, DFS, Binary Search, Topological Sort, Shortest Path, Union Find\\
    数组技巧（强）    & Two Pointers, Sliding Window, Prefix Sum, Prefix Product, Suffix Product, Difference Array\\
    哈希/结构（强）   & Hash Table, Hash Set, Heap, Trie, Monotonic Stack, Monotonic Queue\\
    位运算/数学（强） & Bit Manipulation, Number Theory, Combinatorics\\
    排序/字符串（强） & Counting Sort, Merge Sort, String Matching\\
    \midrule
    通用数据类型（弱） & Array, String, Math, Matrix, Linked List, Stack, Queue, Deque, Tree, Binary Tree, Graph\\
    宽泛范式（弱）     & Recursion, Sorting, Simulation, Brute Force, Regular Expression\\
    \bottomrule
  \end{tabular}
\end{table}
